\clearpage
\begin{landscape}
\begingroup
\fontsize{10}{13}\selectfont
\fancyhead[L]{\small NACHLASS. STÖRUNGEN DER PALLAS}
\renewcommand{\headrulewidth}{0pt}
\pdfbookmark[2]{Allgemeine Störungen, Erste Rechnung, [12]}{pallas-V-fragment12}
\begin{center}[12.]\end{center}
\par\noindent\hspace*{\parindent}[Zur Berechnung der Störungen des zweiten Theils \(\varepsilon\) der Epoche wird die Integration des Ausdrucks \(2r\cos\varphi\mathord{.}\,T\) erfordert; nachdem dieser durch die Interpolationsrechnung hergestellt ist, ergibt seine Integration den Wert von
\[
\begin{aligned}
&\varepsilon-(1-\cos i)\AscendingNode-(1-\cos\varphi)\{\tilde\omega-(1-\cos i)\AscendingNode\}\\
&\qquad=\varepsilon-(1-\cos\varphi)\tilde\omega-\cos\varphi(1-\cos i)\AscendingNode.
\end{aligned}
\]
\par\noindent Die Integration für unser Beispiel:
\begingroup
\newcommand{\VTwelvePoint}{\mathord{.}\mkern5mu}
\newcommand{\VTwelveFive}{\VTwelvePoint\VTwelvePoint\VTwelvePoint\VTwelvePoint\VTwelvePoint}
\begin{center}
\begin{tabular}{@{}l@{\quad}l@{}}
\(\log b\;\VTwelveFive\VTwelveFive\VTwelvePoint0{,}69697\)&\\
\(\log b'\;\VTwelveFive\VTwelveFive1{,}10435_n\)&\\
\multicolumn{1}{r}{\(Q=158^0\,37'\,29''\)}&\\
\(\log\cos Q\;\VTwelveFive\VTwelvePoint\VTwelvePoint9{,}96905_n\)&also das Störungsglied:\\
\(\log q\;\VTwelveFive\VTwelveFive\VTwelvePoint1{,}13530\)&\(+243''{,}02\cos(5M'-2M+158^0\,37'\,29'').\)\\
\(\log\left(5\frac{n'}{n}-2\right)\;\VTwelvePoint\VTwelvePoint\VTwelvePoint8{,}74966_n\)&\\
\cline{1-1}
\multicolumn{1}{r}{\(2{,}38564_n\)}&
\end{tabular}
\end{center}
\endgroup
\par\noindent\hspace*{\parindent}Die gesammten Störungen ergeben sich, wie folgt:
\[
\delta\varepsilon-(1-\cos i)\delta\AscendingNode-(1-\cos\varphi)\{\delta\tilde\omega-(1-\cos i)\delta\AscendingNode\}=
\]
\begingroup
\renewcommand{\arraystretch}{1.15}
\begin{longtable}{@{}l@{}}
\(-21''{,}22\cos(M+349^0\,0'\,41'')\)\\
\(-12{,}47\cos(2M+147\;32\;36)\)\\
\(-2{,}73\cos(3M+212\;6\;43)\)\\[3pt]
\(+2''{,}20\cos(M'-4M+231^0\,22'\,46'')\)\\
\(+5{,}76\cos(M'-3M+290\;31\;20)\)\\
\(+13{,}43\cos(M'-2M+95\;34\;53)\)\\
\(+12{,}20\cos(M'-M+258\;52\;47)\)\\
\SourceLeaf{532}\(-32''{,}95\cos(M'+227^0\,37'\,0'')\)\\
\(-16{,}18\cos(M'+M+37^0\,48'\,32'')\)\\
\(-3{,}40\cos(M'+2M+147\;45\;30)\)\\
\(-1{,}38\cos(M'+3M+266\;33\;49)\)\\[3pt]
\(+1{,}37\cos(2M'-5M+125^0\,11'\,47'')\)\\
\(+1{,}41\cos(2M'-4M+165\;16\;16)\)\\
\(+11{,}78\cos(2M'-3M+18\;6\;21)\)\\
\(+46{,}54\cos(2M'-2M+61\;23\;46)\)\\
\(+196{,}34\cos(2M'-M+223\;43\;43)\)\\
\(-35{,}01\cos(2M'+311\;35\;32)\)\\
\(-6{,}85\cos(2M'+M+73\;15\;56)\)\\
\(-2{,}50\cos(2M'+2M+179\;17\;29)\)\\[3pt]
\(+6{,}34\cos(3M'-4M+274^0\,19'\,19'')\)\\
\(+14{,}42\cos(3M'-3M+329\;2\;43)\)\\
\(+48{,}30\cos(3M'-2M+110\;35\;39)\)\\
\(-166{,}00\cos(3M'-M+227\;50\;23)\)\\
\(-13{,}62\cos(3M'+343\;57\;59)\)\\
\(-3{,}67\cos(3M'+M+91\;40\;13)\)\\
\(-1{,}07\cos(3M'+2M+203\;0\;9)\)\\[3pt]
\(+3{,}02\cos(4M'-5M+170^0\,51'\,10'')\)\\
\(+4{,}94\cos(4M'-4M+247\;47\;37)\)\\
\(+18{,}25\cos(4M'-3M+10\;0\;22)\)\\
\(+52{,}07\cos(4M'-2M+133\;4\;17)\)\\
\(-28{,}81\cos(4M'-M+253\;25\;4)\)\\
\(-5{,}55\cos(4M'+4\;18\;12)\)\\
\(-1{,}64\cos(4M'+M+117\;14\;15)\)\\[3pt]
\(+1{,}30\cos(5M'-6M+71^0\,7'\,44'')\)\\
\(+2{,}26\cos(5M'-5M+172\;30\;27)\)\\
\(+7{,}37\cos(5M'-4M+272\;34\;34)\)\\
\(+16{,}20\cos(5M'-3M+37\;10\;53)\)\\
\(+243{,}02\cos(5M'-2M+158\;37\;29)\)\\
\(-9{,}10\cos(5M'-M+276\;23\;14)\)\\
\(-2{,}33\cos(5M'+30\;8\;31)\)\\[3pt]
\(+1''{,}16\cos(6M'-6M+89^0\,11'\,19'')\)\\
\(+3{,}05\cos(6M'-5M+180\;19\;0)\)\\
\(+6{,}75\cos(6M'-4M+301\;22\;11)\)\\
\(+16{,}06\cos(6M'-3M+63\;27\;50)\)\\
\(-23{,}13\cos(6M'-2M+182\;59\;45)\)\\
\(-3{,}48\cos(6M'-M+300\;6\;0)\)\\
\(-1{,}06\cos(6M'+53\;3\;7)\)\\[3pt]
\(+1{,}35\cos(7M'-6M+96^0\,25'\,40'')\)\\
\(+3{,}03\cos(7M'-5M+205\;49\;26)\)\\
\(+5{,}92\cos(7M'-4M+327\;20\;47)\)\\
\(+22{,}49\cos(7M'-3M+88\;31\;40)\)\\
\(-5{,}85\cos(7M'-2M+206\;57\;35)\)\\
\(-1{,}46\cos(7M'-M+322\;28\;7)\)\\[3pt]
\(+1{,}28\cos(8M'-6M+119^0\,51'\,46'')\)\\
\(+2{,}46\cos(8M'-5M+237\;40\;9)\)\\
\(+4{,}91\cos(8M'-4M+359\;41\;7)\)\\
\(-29{,}31\cos(8M'-3M+121\;55\;2)\)\\
\(-1{,}78\cos(8M'-2M+244\;10\;22)\)\\[3pt]
\(+1{,}24\cos(9M'-6M+137^0\,17'\,43'')\)\\
\(+2{,}19\cos(9M'-5M+258\;13\;31)\)\\
\(+5{,}54\cos(9M'-4M+19\;3\;41)\)\\
\(-3{,}96\cos(9M'-3M+138\;2\;12)\)\\[3pt]
\(+1{,}04\cos(10M'-6M+163^0\,5'\,51'')\)\\
\(+1{,}84\cos(10M'-5M+284\;8\;17)\)\\
\(+14{,}13\cos(10M'-4M+43\;55\;42)\)\\
\(-1{,}20\cos(10M'-3M+162\;54\;33)\)\\[3pt]
\(+1{,}96\cos(11M'-5M+303^0\,37'\,21'')\)\\
\(+3{,}92\cos(11M'-4M+241\;39\;36)\)
\end{longtable}
\begin{center}
\begin{tabular}{@{}lrl@{}}
Säcularänderung: & \(-15''{,}278\)&jährlich\\
&\(-0{,}04183\)&täglich
\end{tabular}
\end{center}
\endgroup
aus welchem Ausdruck man \(\delta\varepsilon\) erhält, wenn man \((1-\cos\varphi)\delta\tilde\omega\) und \(\cos\varphi(1-\cos i)\delta\AscendingNode\) hinzufügt. Den säcularen Theil von \(\delta\varepsilon\) kann man bei Seite lassen, wenn man \(n\) dementsprechend ändert.]
\clearpage
\endgroup
\end{landscape}
